\section{The search budget}
\label{sec:budget}

\subsection{Method}

One bump hunt is one scanned mass spectrum. For a spectrum with constant fractional mass resolution
$\sigma_M = r\,M$, the number of effectively independent resolution elements across a scanned
window $[M_{\mathrm{lo}}, M_{\mathrm{hi}}]$ is
\begin{equation}
  n_s \;=\; \frac{1}{r}\,\ln\!\frac{M_{\mathrm{hi}}}{M_{\mathrm{lo}}}\,,
  \label{eq:ns}
\end{equation}
summed over disjoint segments where a scan is not contiguous. The program's trials factor is
$\Nsig = \sum_s n_s$ and the discovery bar follows from Eq.~\eqref{eq:lee}. This is the
fixed-resolution-element form of the Gross--Vitells estimate~\cite{grossvitells}; the up-crossing
refinement adds a mild $Z$-dependence that is immaterial at the precision quoted here.

The fractional resolution $r$ is the only genuine physics input, so every headline is quoted with
a band obtained by varying $r$ by a factor two in each direction, which moves $\Nsig$ by a factor
two the other way.

\paragraph{The window convention, and why it matters.} Which range counts as ``scanned'' is the
one choice that can quietly bias the whole exercise. We count every spectrum over its
\textbf{published-search scan window}, meaning the mass range a real, published ATLAS search
family scans on that axis, extended where that family's published signal grid demonstrably
prepares a wider scan. Each window carries a per-channel source note
(Appendix~\ref{app:windows}).

The tempting alternative --- taking the range over which signal samples happen to have been
generated --- is wrong twice over. It understates sparsely sampled channels, since a search does
not become cheaper because few samples were made for it, and it ties a supposedly model-space
quantity to one collaboration's production choices. With the published-window convention the
budget is a statement about the searches the field has defined.

\paragraph{Scope.} Only invariant- and transverse-mass bump hunts enter the sum. The dedicated
non-bump programs, meaning displaced HNL, $Z_d$, RPV and ALP searches, MET$+$jet, $dE/dx$ monopole
and off-shell EFT, carry their own look-elsewhere effect and are not mass scans. They are excluded
rather than summed.

\paragraph{What counts as one spectrum.} Models are grouped by the mass axis in which they peak,
and two properties that a real analysis must select on are treated as part of that axis rather than
as cuts within it: \textbf{lepton flavour} and \textbf{$b$-tag content}. So $m(ee)$ and
$m(\mu\mu)$, and likewise $m(eb)$ and $m(ej)$, are separate spectra, each with its own resolution
and scan window. The justification is physical rather than bookkeeping: an electron pair is
measured in the electromagnetic calorimeter ($r \approx 0.015$, nearly mass-independent) while a
muon pair is measured by track sagitta ($r$ rising to $\sim$15\,\% at $3\TeV$, entered as a
window-averaged $r \approx 0.05$), and the two are scanned by separate analyses with separate
triggers. At low mass the ordering reverses, which is why the dark-photon electron channel starts
at $1\GeV$ and the muon channel at $0.3\GeV$. Because flavour is an axis here, it is \emph{not}
also counted in the event-selection multiplicities below; doing both would double-count it.

Conversely, decay modes that resolve onto the \emph{same} mass axis are merged: $X \to HH$ observed
in $4b$ is not a different spectrum from $X \to HH$, it is a different event selection of it, and it
is counted as such in Table~\ref{tab:granularity}.

\subsection{Result}

Table~\ref{tab:budget} gives the per-spectrum breakdown and Table~\ref{tab:budget_summary} the
totals. Covering every bump channel that public BSM models motivate costs
\begin{equation*}
  \Nsig = 3\,694 \quad\text{over 46 spectra}
  \qquad\Longrightarrow\qquad
  \Zloc = 6.44 \ \text{for a } 5\sigma \text{ global discovery.}
\end{equation*}
Equivalently, a local $5\sigma$ anywhere in that program is worth about $3\sigma$ globally.

\begin{table}[htbp]
\centering
\small
\caption{Trials per spectrum. $r$ is the assumed fractional mass resolution, the window is the
published-search scan range (Appendix~\ref{app:windows}) and $n_s$ follows from
Eq.~\eqref{eq:ns}. ``envelope'' is the reference count over the full kinematically accessible
range (analyzable floor to $\sqrt{s}$), ``\#mod'' the number of public model classes that predict a
peak in the spectrum, and ``\#sel'' the number of distinct published event selections that scan it
(Table~\ref{tab:granularity}).}
\label{tab:budget}
\begin{tabular}{rlcrrrrrr}
\toprule
\# & spectrum & $r$ & window [GeV] & $n_s$ & envelope & \#mod & \#sel & $n_s\times$sel \\
\midrule
1 & $m(\gamma\gamma)$ & 0.01 & 60--110 $+$ 150--5000 & 411 & 583 & 5 & 2 & 822\\
2 & $m(\mu\mu)$ ($Z_d$) & 0.02 & 0.3--400 & 360 & 536 & 2 & 1 & 360\\
3 & $m(ee)$ & 0.015 & 150--8000 & 265 & 321 & 7 & 1 & 265\\
4 & $m(ee)$ ($Z_d$) & 0.025 & 1--400 & 240 & 381 & 1 & 1 & 240\\
5 & $m(e\gamma)$ & 0.015 & 200--5000 & 215 & 328 & 1 & 1 & 215\\
6 & $m(e\mu)$ LFV & 0.03 & 100--8000 & 146 & 234 & 4 & 1 & 146\\
7 & $m(e^\pm e^\pm)$ SS & 0.015 & 200--1300 & 125 & 454 & 2 & 1 & 125\\
8 & $m(V\gamma)$ & 0.03 & 220--6000 & 110 & 141 & 3 & 2 & 220\\
9 & $m(ej)$ & 0.03 & 200--5000 & 107 & 234 & 3 & 4 & 429\\
10 & $m(\mu\gamma)$ & 0.03 & 200--5000 & 107 & 164 & 1 & 1 & 107\\
11 & $m(\mu j)$ & 0.04 & 200--5000 & 80 & 176 & 3 & 4 & 322\\
12 & $m(eZ)$ & 0.03 & 100--1100 & 80 & 164 & 1 & 1 & 80\\
13 & $m(\mu\mu)$ & 0.05 & 150--8000 & 80 & 96 & 7 & 1 & 80\\
14 & $m(cb)$ dijet & 0.05 & 20--1000 & 78 & 136 & 1 & 2 & 156\\
15 & $m(e^\pm\mu^\pm)$ SS & 0.025 & 200--1300 & 75 & 272 & 2 & 1 & 75\\
16 & $m(jj)$ & 0.05 & 200--8000 & 74 & 84 & 8 & 2 & 148\\
17 & $m(j\gamma)$ & 0.04 & 500--7000 & 66 & 113 & 2 & 1 & 66\\
18 & $m(\mu Z)$ & 0.04 & 100--1100 & 60 & 123 & 1 & 1 & 60\\
19 & $m(eb)$ & 0.04 & 200--2000 & 58 & 170 & 3 & 2 & 115\\
20 & $m(eejj)$ & 0.05 & 400--7000 & 57 & 76 & 2 & 1 & 57\\
21 & $m(VV)$ & 0.06 & 200--6000 & 57 & 74 & 12 & 9 & 510\\
22 & $m(bj)$ & 0.06 & 200--6000 & 57 & 70 & 1 & 2 & 113\\
23 & $m(Vh)$ & 0.06 & 200--5000 & 54 & 70 & 2 & 6 & 322\\
24 & $m(\mu^\pm\mu^\pm)$ SS & 0.04 & 200--1300 & 47 & 170 & 2 & 1 & 47\\
25 & $m(\mu b)$ & 0.05 & 200--2000 & 46 & 136 & 3 & 2 & 92\\
26 & multilepton & 0.12 & 50--10000 & 44 & 62 & 3 & 4 & 177\\
27 & $m(3j)$ & 0.05 & 200--1800 & 44 & 84 & 2 & 1 & 44\\
28 & $m(tb)$ & 0.08 & 180--6000 & 44 & 54 & 3 & 3 & 132\\
29 & $m(\tau\tau)$ & 0.12 & 35--6000 & 43 & 50 & 4 & 3 & 129\\
30 & $m(\mu\mu jj)$ & 0.07 & 400--7000 & 41 & 54 & 2 & 1 & 41\\
31 & $m(bb)$ & 0.1 & 15--62 $+$ 450--6000 & 40 & 68 & 4 & 3 & 120\\
32 & $m(HH)$ & 0.08 & 250--6000 & 40 & 50 & 4 & 6 & 238\\
33 & $m_{\mathrm{T}}(e\nu)$ & 0.1 & 150--7000 & 38 & 60 & 2 & 1 & 38\\
34 & $m(t\bar t)$ & 0.08 & 350--6000 & 36 & 46 & 9 & 4 & 142\\
35 & $m(\tau j)$ & 0.1 & 200--5000 & 32 & 56 & 3 & 2 & 64\\
36 & $m(e\tau)$ LFV & 0.12 & 100--3000 & 28 & 50 & 4 & 1 & 28\\
37 & $m(\mu\tau)$ LFV & 0.12 & 100--3000 & 28 & 50 & 4 & 1 & 28\\
38 & $m_{\mathrm{T}}(\tau\nu)$ & 0.12 & 200--5000 & 27 & 47 & 1 & 1 & 27\\
39 & $m_{\mathrm{T}}(\mu\nu)$ & 0.15 & 150--7000 & 26 & 40 & 2 & 1 & 26\\
40 & $m(tW)$ & 0.08 & 500--3000 & 22 & 50 & 1 & 2 & 45\\
41 & $m(\mathrm{multi})$ & 0.06 & 3000--10500 & 21 & 64 & 2 & 1 & 21\\
42 & $m(t\bar t)/m(jj)$ (sgluon) & 0.08 & 400--2000 & 20 & 53 & 2 & 2 & 40\\
43 & $m(\tau b)$ & 0.12 & 200--2000 & 19 & 47 & 3 & 2 & 38\\
44 & $m(t\bar tZ)/m(Zt)$ & 0.08 & 1000--4000 & 17 & 43 & 1 & 2 & 35\\
45 & $m(Wb)$ & 0.08 & 800--3000 & 16 & 61 & 1 & 1 & 16\\
46 & $m(Ht)$ & 0.08 & 1000--3000 & 14 & 48 & 1 & 1 & 14\\
\midrule
\multicolumn{4}{l}{\textbf{total (46 spectra)}} & \textbf{3\,694} & 6\,442 & & \textbf{94} & \textbf{6\,616}\\
\bottomrule
\end{tabular}
\end{table}

\begin{table}[htbp]
\centering
\caption{Search-budget summary. The band comes from varying the fractional resolution $r$ by a
factor two in each direction.}
\label{tab:budget_summary}
\begin{tabular}{lrrlrl}
\toprule
estimate & spectra & $\Nsig$ & band & $\Zloc$ & band \\
\midrule
inclusive (one spectrum per observable) & 46 & 3694 & 1847--7388 & \textbf{6.44} & 6.33--6.54 \\
published event selections & 94 & 6616 & 3308--13232 & \textbf{6.53} & 6.42--6.63 \\
kinematic envelope (reference bound) & 46 & 6442 & 3221--12884 & 6.52 & 6.42--6.63 \\
\bottomrule
\end{tabular}
\end{table}

\begin{figure}[htbp]
\centering
\includegraphics[width=0.92\textwidth]{search_budget.png}
\caption{Trials per spectrum, with the scan window annotated and the colour giving the number of
public model classes that feed the channel. Note the strong hierarchy: the five highest-resolution
spectra ($m(\gamma\gamma)$, the two dark-photon axes, $m(ee)$ and $m(e\gamma)$) carry $40\,\%$ of
the entire budget, simply because $n_s \propto 1/r$.}
\label{fig:budget}
\end{figure}

\subsection{Which models feed which spectrum}

The mapping from model classes to spectra is the other half of the input, and it is worth being able
to see it rather than trust it. Figure~\ref{fig:matrix} draws it in full. Two structural features
matter for the budget. First, the matrix is far from diagonal: popular spectra such as $m(jj)$,
$m(VV)$ and $\mtt$ are motivated by eight to twelve model classes each, which is exactly why the
number of \emph{searches} is so much smaller than the number of \emph{models}. Second, adding a
model to a spectrum that is already scanned costs \emph{nothing} in trials. The entire cost of
breadth sits in opening new columns, never in filling in existing ones.

\begin{figure}[htbp]
\centering
\includegraphics[width=0.95\textwidth]{model_observable_matrix.png}
\caption{Public model classes (rows) against canonical bump spectra (columns). A filled cell means
the model predicts a peak in that spectrum. The matrix is derived from the model map rather than
asserted by hand.}
\label{fig:matrix}
\end{figure}

\subsection{How finely should searches be counted?}

An inclusive count of one spectrum per observable is the lower bound. Real searches slice each
final state further, and each slice scanned for a bump is another look. Published searches split by
$b$-tag category, boosted versus resolved regime, sub-decay mode, $0/1/2$-lepton bin and ambiguous
pairing --- a leptoquark pair search scans all four lepton--jet pairings. Attaching a documented,
publication-anchored selection multiplicity to each observable gives 94 channels and
$\Nsig = 6\,616$, hence $\Zloc = 6.53$. Lepton flavour is deliberately \emph{not} one of the
multipliers, since it is already a mass axis.

\begin{table}[htbp]
\centering
\caption{The budget at increasing granularity.}
\label{tab:granularity}
\begin{tabular}{lrrr}
\toprule
granularity & spectra & $\Nsig$ & $\Zloc$ \\
\midrule
inclusive (one spectrum per observable) & 46 & 3694 & 6.44 \\
real event selections & 94 & 6616 & \textbf{6.53} \\
full ATLAS BSM program (literature, Section~\ref{sec:excess}) & --- & $\sim5\times10^4$ & 6.83 \\
\bottomrule
\end{tabular}
\end{table}

\begin{figure}[htbp]
\centering
\includegraphics[width=0.88\textwidth]{budget_waterfall.png}
\caption{$\Zloc$ required for a $5\sigma$ global discovery as a function of the trials factor, with
the granularity levels of this study placed on the curve. A factor $1.8$ in channel counting moves
the discovery bar by $0.09\sigma$, and reaching the full ATLAS program, a further factor $\sim8$,
adds only another $0.3\sigma$.}
\label{fig:waterfall}
\end{figure}

\subsection{Reading the result}

The headline is robust by construction. $\Nsig$ enters only through $\ln \Nsig$, so a factor two
error in the resolution assumption, or a factor $1.8$ in how finely channels are counted, moves
the discovery bar by at most $\pm 0.1\sigma$. Every granularity level of this study in
Table~\ref{tab:granularity} sits within $0.05\sigma$ of $6.5$.

Two conclusions follow, and they pull in opposite directions.

\textbf{Breadth is cheap.} Covering the entire known BSM resonance space costs about one extra
unit of $\sigma$ relative to a single pre-specified counting experiment, $6.5$ against $5$. There
is no statistical argument for keeping the search program narrow. The marginal cost of scanning
one more spectrum is small, and the cost of adding models to an existing one is exactly zero.

\textbf{Uncountability is expensive.} The argument above works only because $\Nsig$ can be
counted. The moment the number of effective looks becomes a matter of opinion, which is the
natural state of a neural-network-driven scan, the correction one has to defend is no longer $2\ln
\Nsig$ with $\Nsig$ known but whatever a reviewer will accept. Section~\ref{sec:ab} quantifies how
much that is worth paying to avoid.
